\section{Lecciones para el proyecto propio}
\label{sec:lecciones}

\begin{table}[H]
  \centering
  \small
  \caption{Riesgos del proyecto y medidas preventivas}
  \label{tab:riesgos}
  \begin{tabularx}{\textwidth}{@{}p{0.6cm}LL@{}}
    \toprule
    \textbf{N.\textdegree} & \textbf{Riesgo concreto} & \textbf{Medida preventiva} \\
    \midrule
    1 &  \textbf{Sistema de Pre-Sustentaciones:} si el acta se genera desde una réplica atrasada, como las que en GitHub quedaron horas detrás del primario \cite{Warner2018}, sale sin la nota que el tribunal acaba de registrar y se firma con ese error. &En TA-PFC-E1 se define que notas y actas se leen siempre del primario. Una prueba lo verifica: registra una nota y genera el acta enseguida.\\
    2 & & \\
    3 & \textbf{Registro de notas:} si se corta la conexión mientras un docente guarda una nota, el sistema podría no registrarla aunque parezca que sí. & Confirmar el guardado antes de mostrar un mensaje de éxito y probar cortes de conexión para evitar registros perdidos o duplicados. \\
    \bottomrule
  \end{tabularx}
\end{table}
